\documentclass[french]{book}
\usepackage{teach}
\usepackage{verbatim}
\usepackage{amsmath,amsfonts,amssymb}
\usepackage{pgfplots}
%\usepackage{geometry}
%\geometry{top=1cm, bottom=1cm, left=2cm, right=2cm}
\usepackage[utf8]{inputenc}
\usepackage[french]{babel}
\redefinecolor{thm}{title}{bgcolor}{alizarin}
\redefinecolor{prop}{title}{bgcolor}{meatbrown}
\redefinecolor{defn}{title}{bgcolor}{olivine}
\definecolor{alizarin}{rgb}{0.82, 0.1, 0.26}
\definecolor{meatbrown}{rgb}{0.9, 0.72, 0.23}
\definecolor{olivine}{rgb}{0.6, 0.73, 0.45}
\usepackage{pgf,tikz,tkz-tab}

\usepackage{mathrsfs}
\usetikzlibrary{arrows}

\begin{document}

\chapter{Calcul littéral}
\section{Puissances}
\begin{defn}
Soient $a\in\mathbb{R}$ un nombre réel et $n\in \mathbb{N}^{*}$. On définit la puissance de la manière suivante:

\begin{itemize}
\item[$\bullet$]$a^n=\underbrace{a\times a \times ... \times a}_{n \, \rm fois}$
\item[$\bullet$] Si $a \neq 0$, $a^{-n}=\frac{1}{a^n}=\frac{1}{a\times a \times ... \times a}$. 
\item[$\bullet$] Si $a\neq 0$, $a^0=1$.
\end{itemize}
\end{defn}

\underline{Exemple:}
\begin{itemize}
\item[$\bullet$]$2^{-3}=\frac{1}{2^3}=\frac{1}{2\times 2 \times 2}=\frac{1}{8}$
\item[$\bullet$] Le nombre $\frac{1}{5}$ peut s'écrire $5^{-1}$. C'est l'inverse de $5$.

\item[$\bullet$] La décomposition de 4200 en produit de nombre premiers peut s'écrire:
$4200=2\times2\times2\times3\times5\times5\times7=2^3\times3\times5^2\times7$
\end{itemize}

\begin{prop}
\begin{itemize}
\item[$\bullet$] $a^m \times a^n = a^{m+n} $
\item[$\bullet$] $\frac{a^m}{a^n}= a^{m-n} $
\item[$\bullet$] $(a^m)^n=a^{m\times n}$
\item[$\bullet$] $a^m \times b^m = (ab)^m $
\item[$\bullet$] $\frac{a^m}{b^m}=(\frac{a}{b})^m$
\end{itemize}
\end{prop}

\underline{Exemple :}
\begin{itemize}
\item[$\bullet$]$(-3)^4 \times (-3)^2= (-3)^{4+2}=(-3)^6$
\item[$\bullet$]$\frac{2^7}{2^5}=2^{7-5}=2^2$
\item[$\bullet$]$(5^4)^3=5^{3 \times 4}=5^{12}$
\item[$\bullet$]$10^5\times 6^5=(10\times 6)^5=60^5$
\item[$\bullet$]$\frac{(-8)^3}{2^3}=(\frac{-8}{2})^3=(-4)^3$
\end{itemize}

\begin{defn}
Soit $x \in \mathbb{D}^*$. On dit que l'écriture scientifique de $x$ est $a \times 10^n$ où $n \in \mathbb{N}$ et $a\in\mathbb{D}$ tel que $1\leq \vert a \vert < 10$.  \vspace{4mm}
\end{defn}
 
\underline{Exemple :}
L'écriture scientifique de $0,000\,45$ est $4,5 \times 10^{-4}$.

\subsection{La racine carrée}
\begin{defn}
La racine carrée d'un nombre $a\in\mathbb{R}_+$ est l'unique réel positif dont le carré vaut $a$. La racine carrée de $a$ se note $\sqrt{a}$ et $\sqrt{a}\in \mathbb{R}_+$.
\end{defn}

\begin{prop}
Soit $x\in\mathbb{R}$, on a $\sqrt{x^2}=\vert x \vert$.
Soient $a$ et $b$ deux réels positifs. On a $\sqrt{ab}=\sqrt{a}\times \sqrt{b}$, et si $b\neq0$, $\sqrt{\frac{a}{b}}=\frac{\sqrt{a}}{\sqrt{b}}$
\end{prop}

\begin{demo}
Soient $a,b \in \mathbb{R}_+$, on sait d'une part que:
$$\sqrt{ab} \times \sqrt{ab} = \vert ab \vert = ab $$
d'autre part on a aussi que:
$$\sqrt{a} \sqrt{b} \times \sqrt{a} \sqrt{b}= \sqrt{a} \sqrt{a} \times \sqrt{b} \sqrt{b}= a \times b = ab$$
par unicité de la racine carré on a:
$$\sqrt{ab}=\sqrt{a} \sqrt{b}$$
\end{demo}

\underline{Remarque :}
si $a$ et $b$ sont non nuls alors $\sqrt{a+b} \neq \sqrt{a} + \sqrt{b}$

\begin{thm}
Soit $a,b \in \mathbb{R}^{*}_{+}$ alors on a $\sqrt{a+b} < \sqrt{a} + \sqrt{b}$.
\end{thm}

\begin{demo}
La distance la plus courte entre deux points $B$ et $C$ du plan est toujours la longueur du segment $[BC]$,  ainsi tout autre chemin allant de $B$ à $C$ a une longueur nécessairement plus grande que la longueur du segment $[BC]$, par suite on place $A$ un point du plan tel que $CAB$ soit un triangle rectangle non applati. \\

On pose $CA=\sqrt{a} \; (>0)$ et $AB=\sqrt{b} \; (>0)$. Par le théorème de Pythagore on obtient que $BC=\sqrt{a+b}$ or $BC$ étant la distance la plus courte elle sera strictement plus petite que la longueur du chemin allant de $B$ à $C$ formé par $CA$ et $AB$ qui est égale à $\sqrt{a}+\sqrt{b}$.  D'où le résultat.

\definecolor{qqwuqq}{rgb}{0,0.39215686274509803,0}
\definecolor{ududff}{rgb}{0.30196078431372547,0.30196078431372547,1}
\begin{center}
\begin{tikzpicture}[scale=1.1,line cap=round,line join=round,>=triangle 45,x=1cm,y=1cm]
\clip(-9.008933262993455,-1.975143952375375) rectangle (-4.622531938487147,2.377750846179176);
\draw[line width=1pt,color=qqwuqq,fill=qqwuqq,fill opacity=0.10000000149011612] (-4.967049946429461,-1.5227455665703808) -- (-5.10229116463362,-1.5239460261329616) -- (-5.101090705071039,-1.6591872443371203) -- (-4.9658494868668805,-1.6579867847745395) -- cycle; 
\draw [line width=1pt] (-4.997255209576774,1.88011506410132)-- (-4.9658494868668805,-1.6579867847745395);
\draw [line width=1pt] (-4.9658494868668805,-1.6579867847745395)-- (-8.503951335742741,-1.6893925074844327);
\draw [line width=1pt] (-8.503951335742741,-1.6893925074844327)-- (-4.997255209576774,1.88011506410132);
\begin{scriptsize}
\draw [fill=ududff] (-4.9658494868668805,-1.6579867847745395) circle (1pt);
\draw[color=ududff] (-4.704804104666275,-1.654805604039167) node {$A$};
\draw [fill=ududff] (-4.997255209576774,1.88011506410132) circle (1pt);
\draw[color=ududff] (-4.947686687832818,2.0175303885707385) node {$B$};
\draw [fill=ududff] (-8.503951335742741,-1.6893925074844327) circle (1pt);
\draw[color=ududff] (-8.556266847237717,-1.5272939376291008) node {$C$};
\draw[color=qqwuqq] (-5.434587937255469,-1.342402021334505) node {$\alpha = 90\textrm{\degre}$};
\draw[color=black] (-4.80043604897215,0.1558600589837726) node {$\sqrt{b}$};
\draw[color=black] (-6.6818453510097315,-1.85823172704492332) node {$\sqrt{a}$};
\draw[color=black] (-7.1371017111879,0.29612289203484543) node {$\sqrt{a+b}$};
\end{scriptsize}
\end{tikzpicture}
\end{center}

\end{demo}

\definecolor{qqwuqq}{rgb}{0,0.39215686274509803,0}
\definecolor{ududff}{rgb}{0.30196078431372547,0.30196078431372547,1}
\begin{comment}
\begin{tikzpicture}[line cap=round,line join=round,>=triangle 45,x=1cm,y=1cm]
\clip(-6.766033237937672,-5.384367780089153) rectangle (4.358412086886164,5.644271381623621);
\draw[line width=2pt,color=qqwuqq,fill=qqwuqq,fill opacity=0.10000000149011612] (3.0971469769010596,0.33546194894073955) -- (2.713832503981273,0.33229406073479095) -- (2.717000392187222,-0.05102041218499551) -- (3.1003148651070083,-0.04785252397904695) -- cycle; 
\draw [line width=2pt] (-1.2726819131655425,-0.08399299322096865)-- (3.0641743958650864,4.325144254293504);
\draw [line width=2pt] (3.1003148651070083,-0.04785252397904695)-- (-1.2726819131655425,-0.08399299322096865);
\draw [line width=2pt] (3.1003148651070083,-0.04785252397904695)-- (3.0641743958650864,4.325144254293504);
\begin{scriptsize}
\draw [fill=ududff] (-1.2726819131655425,-0.08399299322096865) circle (2.5pt);
\draw[color=ududff] (-1.128120036197855,0.3045170511296897) node {$A$};
\draw [fill=ududff] (3.1003148651070083,-0.04785252397904695) circle (2.5pt);
\draw[color=ududff] (3.3894386190423833,-0.25566022212009676) node {$B$};
\draw [fill=ududff] (3.0641743958650864,4.325144254293504) circle (2.5pt);
\draw[color=ududff] (3.281017211316618,4.713654298644138) node {$C$};
\draw[color=qqwuqq] (2.529261345792594,0.8104836205165936) node {$\alpha = 90$};
\draw[color=black] (0.20907732575325563,2.3090856748869137) node {$\sqrt{a+b}$};
\draw[color=black] (0.7511843643820844,-0.31951975287817505) node {$\sqrt{a}$};
\draw[color=black] (3.6424219037358365,2.1657512170886575) node {$\sqrt{b}$};
\end{scriptsize}
\end{tikzpicture}
\end{comment}

\newpage
\section{Expressions algébriques}

\subsection{Développer, Factoriser}
\begin{prop}
Soient $a,b,c,d,k \in \mathbb{R}$ des nombres réels, alors on a: 
\begin{itemize}
\item[$\bullet$] $k(a+b)=ka+kb$
\item[$\bullet$] $(a+b)(c+d)=ac+ad+bc+bd$
\end{itemize}
\end{prop}

\underline{Exemple :}

$$2(x+3)=2x+2\times3=2x+6$$
$$(2+x)(4-x)=2\times4-2\times x + x \times 4 - x \times x$$
On calcule en ordonnant nombres puis inconnues dans cet ordre
$$(2+x)(4-x)=8-2x+4x-x^2$$
on termine en réduisant et en réordonnant les termes, dans l'ordre décroissant des exposants de l'inconnue.
$$(2+x)(4-x)=-x^2+2x+8$$

\subsection{Identités remarquables}

Voici ici trois expressions algébriques à connaitre parfaitement.

\begin{prop}
Soient $a,b \in \mathbb{R}$ des nombres réels, alors on a: 
\begin{itemize}
\item[$\bullet$] $(a+b)^2=a^2+2ab+b^2$
\item[$\bullet$] $(a-b)^2=a^2-2ab+b^2$
\item[$\bullet$] $(a-b)(a+b)=a^2-b^2$
\end{itemize}
\end{prop}


\definecolor{qqzzcc}{rgb}{0,0.6,0.8}
\definecolor{ffzzqq}{rgb}{1,0.6,0}
\definecolor{yqqqqq}{rgb}{0.5019607843137255,0,0}
\definecolor{qqwuqq}{rgb}{0,0.39215686274509803,0}
\definecolor{ttqqqq}{rgb}{0.2,0,0}
\begin{center}

\begin{tikzpicture}[scale=1.8,line cap=round,line join=round,>=triangle 45,x=1cm,y=1cm]
\clip(-0.7,-0.9) rectangle (5,4.2);
\fill[line width=2pt,color=white] (-0.38794762556557366,-0.8496045569773465) -- (4.457924829196802,-0.8361728794074994) -- (4.444493151626955,4.009699575354875) -- (-0.40137930313542025,3.996267897785029) -- cycle;
\fill[line width=2pt,color=qqwuqq,fill=qqwuqq,fill opacity=0.7] (-0.40137930313542025,3.996267897785029) -- (2.9901382143185278,4.005668427863396) -- (2.999538744396895,0.6141509104094485) -- (-0.3919787730570524,0.6047503803310805) -- cycle;
\fill[line width=2pt,color=yqqqqq,fill=yqqqqq,fill opacity=0.7] (-0.3919787730570524,0.6047503803310805) -- (2.999538744396895,0.6141509104094485) -- (3.0035698918883735,-0.8402040268989782) -- (-0.38794762556557366,-0.8496045569773465) -- cycle;
\fill[line width=2pt,color=ffzzqq,fill=ffzzqq,fill opacity=0.6] (2.999538744396895,0.6141509104094485) -- (4.453893681705324,0.6181820579009272) -- (4.457924829196802,-0.8361728794074994) -- (3.0035698918883735,-0.8402040268989782) -- cycle;
\fill[line width=2pt,color=qqzzcc,fill=qqzzcc,fill opacity=0.6] (2.999538744396895,0.6141509104094485) -- (2.9901382143185278,4.005668427863396) -- (4.444493151626955,4.009699575354875) -- (4.453893681705324,0.6181820579009272) -- cycle;
\draw [line width=2pt] (-0.40137930313542025,3.996267897785029)-- (-0.3919787730570524,0.6047503803310805);
\draw [line width=2pt] (-0.3412131460999479,2.3006628781389677) -- (-0.45214493009252477,2.300355399977141);
\draw [line width=2pt,color=qqwuqq] (-0.40137930313542025,3.996267897785029)-- (2.9901382143185278,4.005668427863396);
\draw [line width=2pt,color=qqwuqq] (1.2942257165106406,4.0564340548205005) -- (1.2945331946724672,3.945502270827924);
\draw [line width=2pt,color=qqwuqq] (2.9901382143185278,4.005668427863396)-- (2.999538744396895,0.6141509104094485);
\draw [line width=2pt,color=qqwuqq] (3.0503043713539997,2.310063408217336) -- (2.9393725873614236,2.309755930055509);
\draw [line width=2pt,color=qqwuqq] (2.999538744396895,0.6141509104094485)-- (-0.3919787730570524,0.6047503803310805);
\draw [line width=2pt,color=qqwuqq] (1.303933724750835,0.5539847533739759) -- (1.303626246589008,0.6649165373665527);
\draw [line width=2pt,color=qqwuqq] (-0.3919787730570524,0.6047503803310805)-- (-0.40137930313542025,3.996267897785029);
\draw [line width=2pt,color=yqqqqq] (-0.3919787730570524,0.6047503803310805)-- (2.999538744396895,0.6141509104094485);
\draw [line width=2pt,color=yqqqqq] (2.999538744396895,0.6141509104094485)-- (3.0035698918883735,-0.8402040268989782);
\draw [line width=2pt,color=yqqqqq] (3.0569561521885427,-0.08976203083206524) -- (2.9460243681959666,-0.09006950899389195);
\draw [line width=2pt,color=yqqqqq] (3.0570842680893033,-0.13598360749563845) -- (2.946152484096727,-0.13629108565746514);
\draw [line width=2pt,color=yqqqqq] (3.0035698918883735,-0.8402040268989782)-- (-0.38794762556557366,-0.8496045569773465);
\draw [line width=2pt,color=yqqqqq] (1.3079648722423132,-0.9003701839344501) -- (1.3076573940804865,-0.7894383999418744);
\draw [line width=2pt,color=yqqqqq] (-0.38794762556557366,-0.8496045569773465)-- (-0.3919787730570524,0.6047503803310805);
\draw [line width=2pt,color=yqqqqq] (-0.4453650333572209,-0.14569161573583325) -- (-0.3344332493646441,-0.14538413757400656);
\draw [line width=2pt,color=yqqqqq] (-0.44549314925798206,-0.09947003907226006) -- (-0.3345613652654052,-0.09916256091043334);
\draw [line width=2pt,color=ffzzqq] (2.999538744396895,0.6141509104094485)-- (4.453893681705324,0.6181820579009272);
\draw [line width=2pt,color=ffzzqq] (3.7034516856384094,0.6715683182010961) -- (3.7037591638002363,0.5606365342085193);
\draw [line width=2pt,color=ffzzqq] (3.749673262301983,0.6716964341018566) -- (3.7499807404638097,0.5607646501092799);
\draw [line width=2pt,color=ffzzqq] (4.453893681705324,0.6181820579009272)-- (4.457924829196802,-0.8361728794074994);
\draw [line width=2pt,color=ffzzqq] (4.511311089496972,-0.08573088334058582) -- (4.400379305504395,-0.08603836150241252);
\draw [line width=2pt,color=ffzzqq] (4.511439205397732,-0.13195246000415903) -- (4.400507421405156,-0.13225993816598572);
\draw [line width=2pt,color=ffzzqq] (4.457924829196802,-0.8361728794074994)-- (3.0035698918883735,-0.8402040268989782);
\draw [line width=2pt,color=ffzzqq] (3.754011887955288,-0.8935902871991455) -- (3.7537044097934613,-0.7826585032065698);
\draw [line width=2pt,color=ffzzqq] (3.7077903112917148,-0.8937184030999061) -- (3.707482833129888,-0.7827866191073304);
\draw [line width=2pt,color=ffzzqq] (3.0035698918883735,-0.8402040268989782)-- (2.999538744396895,0.6141509104094485);
\draw [line width=2pt,color=qqzzcc] (2.999538744396895,0.6141509104094485)-- (2.9901382143185278,4.005668427863396);
\draw [line width=2pt,color=qqzzcc] (2.9901382143185278,4.005668427863396)-- (4.444493151626955,4.009699575354875);
\draw [line width=2pt,color=qqzzcc] (3.6940511555600413,4.063085835655043) -- (3.694358633721868,3.952154051662467);
\draw [line width=2pt,color=qqzzcc] (3.7402727322236147,4.063213951555804) -- (3.7405802103854415,3.9522821675632276);
\draw [line width=2pt,color=qqzzcc] (4.444493151626955,4.009699575354875)-- (4.453893681705324,0.6181820579009272);
\draw [line width=2pt,color=qqzzcc] (4.504659308662427,2.314094555708814) -- (4.393727524669851,2.3137870775469875);
\draw [line width=2pt,color=qqzzcc] (4.453893681705324,0.6181820579009272)-- (2.999538744396895,0.6141509104094485);
\draw [line width=2pt] (-0.3919787730570524,0.6047503803310805)-- (-0.38794762556557366,-0.8496045569773465);
\draw [line width=2pt] (-0.40137930313542025,3.996267897785029)-- (2.9901382143185278,4.005668427863396);
\draw [line width=2pt] (2.9901382143185278,4.005668427863396)-- (4.444493151626955,4.009699575354875);
\draw [line width=2pt] (-0.3919787730570524,0.6047503803310805)-- (2.999538744396895,0.6141509104094485);
\draw [line width=2pt] (2.999538744396895,0.6141509104094485)-- (2.9901382143185278,4.005668427863396);
\draw [line width=2pt] (2.999538744396895,0.6141509104094485)-- (4.453893681705324,0.6181820579009272);
\draw [line width=2pt] (2.999538744396895,0.6141509104094485)-- (3.0035698918883735,-0.8402040268989782);
\draw [line width=2pt] (4.453893681705324,0.6181820579009272)-- (4.457924829196802,-0.8361728794074994);
\draw [line width=2pt] (3.0035698918883735,-0.8402040268989782)-- (4.457924829196802,-0.8361728794074994);
\draw [line width=2pt] (4.444493151626955,4.009699575354875)-- (4.453893681705324,0.6181820579009272);
\draw [line width=2pt] (4.453893681705324,0.6181820579009272)-- (4.457924829196802,-0.8361728794074994);
\draw [line width=2pt] (-0.38794762556557366,-0.8496045569773465)-- (3.0035698918883735,-0.8402040268989782);
\draw [line width=2pt] (4.457924829196802,-0.8361728794074994)-- (3.0035698918883735,-0.8402040268989782);
\begin{scriptsize}
\draw[color=black] (-0.59935765995896,2.382056519208632) node {$a$};
\draw[color=qqwuqq] (1.323467315492491,2.419033922582698) node {$A_1$};
\draw[color=yqqqqq] (1.3142229646489745,0.024747054111914135) node {$A_2$};
\draw[color=yqqqqq] (-0.5716246074284103,-0.06769645432325125) node {$b$};
\draw[color=ffzzqq] (3.699265482276255,-0.02147470010566857) node {$A_4$};
\draw[color=qqzzcc] (3.7177541839632884,2.437522624269731) node {$A_3$};
\end{scriptsize}
\end{tikzpicture}
\end{center}

\newpage


\section{Équations, inéquations}
\subsection{Équations}
\begin{defn}
Une équation $E$ est une égalité mathématique qui peut faire intervenir des variables et qui n'est pas nécessairement vraie. \\

\underline{Résoudre une équation}, c'est trouver \underline{toutes} les valeurs des variables possibles telles que l'égalité devient vraie. Si l'on met \underline{toutes} ces valeurs dans un même ensemble on obtient alors \underline{l'ensemble des solutions} associé à l'équation $E$, que l'on notera $S$. 
\end{defn}

\subsubsection{Équation du premier degré}
Ici on se donne une variable réelle que l'on note $x$, une équation du premier degré est une égalité du type: 
$$(E) \; ax+b=0 $$ 
où $a\in \mathbb{R}^{*}$ et $b\in \mathbb{R}$.

Notre but est de résoudre ce type d'équation, dans une équation du premier degré on peut décrire directement l'ensemble des solutions $S$.

$$S=\{ - \frac{b}{a}\}$$
On remarque que ce type d'équation admet toujours une \underline{unique} solution.\\


\underline{Exemple}: On souhaite résoudre l'équation suivante, $(E) \; 4x+3=0$, cette équation n'est pas vraie par exemple pour $x=0$ ou encore pour $x=2$.

Il nous faut donc trouver l'ensemble des solutions $S$ pour savoir pour quelle valeur de $x$ l'équation $E$ est vraie. Ici $S=\{-\frac{3}{4}\}$, donc l'équation $E$ est vraie si et seulement si $x\in S$ ou encore ici, si $x=-\frac{3}{4}$.

\subsubsection{Équation produit}
On peut être confronté aussi à résoudre des équations du type $(E) \; (4x+3)(2x+1)=0$ appelé encore équation produit, pour cela on utilise un théorème déjà vu au collège

\begin{thm}
Un produit de facteur est nul \underline{si et seulement si} l'un des facteurs \underline{(au moins)} est nul. Autrement dit:

$$A \times B = 0 \Longleftrightarrow A=0 \;\; \textrm{ou} \;\; B=0$$

\end{thm}

Ainsi résoudre une équation produit revient à résoudre plusieurs équations.
$$(E) \; (4x+3)(2x+1)=0$$
$$(E) \; (4x+3)=0 \;\;  \textrm{ou} \;\; (2x+1)=0$$
$$(E) \; 4x=-3 \;\;  \textrm{ou} \;\; 2x=-1$$
$$(E) \; x=-\frac{3}{4} \;\;  \textrm{ou} \;\; x=-\frac{1}{2}$$
$$ S = \{  -\frac{3}{4} ; -\frac{1}{2} \} $$

\subsection{Inéquations}
\begin{defn}
Une inéquation $I$ est une inégalité mathématique qui peut faire intervenir des variables et qui n'est pas nécessairement vraie. \\

\underline{Résoudre une inéquation}, c'est trouver \underline{toutes} les valeurs des variables possibles tel que l'inégalité soit vraie. Si l'on met \underline{toutes} ces valeurs dans un même ensemble on obtient alors \underline{l'ensemble des solutions} associé à l'inéquation $I$, que l'on notera aussi $S$. 
\end{defn}

\subsubsection{Inéquation du premier degré}

Une différence notable entre la résolution d'inéquation et la résolution d'équations se trouve dans le  \underline{signe des nombres}, qui doit être prise en compte dans les inéquations contrairement au équations donc il n'y aura pas de méthode générale ici.\\


\underline{Exemple}: Résolvons l'inéquation suivante:
$$(I) \; -3x+2\leq 0$$
$$(I) \; -3x\leq -2$$
$$(I) \; x \geq \frac{2}{3}$$

Ainsi on a $S=[\frac{2}{3};+\infty]$. On remarque le changement de signe entre la deuxième ligne et la troisième ligne qui est dû au fait que $-3$ soit négatif.

\subsubsection{Tableau de signe}
De manière analogue au tableau de variation on peut dresser le tableau de signe d'une expression mathématique ce qui nous permettra de connaitre simplement le signe d'expressions parfois plus complexe.

Voici un tableau de signe en reprenant l'exemple précédant:
\begin{center}
\begin{tikzpicture}
   \tkzTabInit{$x$ / 1 , $-3x+2$ / 1}{$-\infty$, $0$, $\frac{2}{3}$, $+\infty$}
   \tkzTabLine{, , +, , z, -, }
\end{tikzpicture}
\end{center}

L'intérêt de cet outil réside dans la résolution d'inéquations plus complexe telles que 

$$(I) \;  (-3x+2)(4x+3)\geq 0$$

Pour cela on rappelle quelques règles mnémotechniques: $"+ \times + donne +"$, $"- \times - donne +"$, $"- \times + donne -"$ et $"+ \times - donne -"$.  

\begin{center}
\begin{tikzpicture}
   \tkzTabInit{$x$ / 1 , $-3x+2$ / 1, $4x+3$ /1 , $(-3x+2)(4x+3)$ /1}{$-\infty$, -$\frac{3}{4}$ ,$0$, $\frac{2}{3}$, $+\infty$}
   \tkzTabLine{,+, , , +, , z , -, }
   \tkzTabLine{,-,z, , +, ,  , +, }
   \tkzTabLine{,-,z, , +, , z , -, }
\end{tikzpicture}
\end{center}

Ainsi l'ensemble des solutions est 
$ S= [-\frac{3}{4}; \frac{2}{3}] $.


\end{document}
